\section{Every sufficiently large bound in the corrected lattice}
\label{sec:projective}

The key distinction is now between qualitative existence and the uniform estimates required by the source. We can prove the former in a general projective setting.

\subsection{Statement}

Let $K$ be imaginary quadratic, $\OO=\OO_K$, and assume that some $\eta\in\OO$ satisfies
\begin{equation}\label{eq:traceone}
 \eta+\overline\eta=1.
\end{equation}
For the source, $K=\Q(\sqrt{-d})$ with $d\equiv7\pmod8$, and $\eta=(1+\sqrt{-d})/2$ works. Let $V$ have positive dimension over $K$, and let $h$ be positive definite Hermitian, conjugate-linear in the first variable. For $v\in V$, define the rank-one operator
\[
 \Theta_v(z)=v h(v,z).
\]
For $x,y\in V$, also put $\Theta_{x,y}(z)=x h(y,z)$.

\begin{theorem}[Projective Hermitian eventual frames]\label{thm:projective}
Let $A$ be a full projective $\OO$-lattice in $V$ with $A^\#=A$. Let $S\subset A$ be finite and $R\in\Z_{\ge0}$. There are a fixed finite family $T\subset A$ containing $S$ and an integer $M_0$ such that for every integer $M\ge M_0$ there are integers $m_v\ge R$ satisfying
\begin{equation}\label{eq:projective-frame}
 \sum_{v\in T}m_v\Theta_v=M I_V.
\end{equation}
In orthonormal complex coordinates this is an ordinary tight-frame identity. Neither freeness of $A$ nor a principal fractional ideal is assumed.
\end{theorem}

\subsection{Step 1: an integral signed resolution of the identity}

\begin{lemma}\label{lem:signed}
Under the hypotheses of Theorem~\ref{thm:projective}, the identity is an integer linear combination of finitely many $\Theta_v$ with $v\in A$.
\end{lemma}

\begin{proof}
The dual-basis lemma for a finite projective module supplies $x_i\in A$ and $\OO$-linear functionals $\ell_i:A\to\OO$ such that
\[
 z=\sum_i x_i\ell_i(z).
\]
Each functional extends to $V$ and, by nondegeneracy, has the form $\ell_i(z)=h(y_i,z)$ for some $y_i\in V$. Since $\ell_i(A)\subseteq\OO$, self-duality implies $y_i\in A$. Thus
\[
 \sum_i\Theta_{x_i,y_i}=I,
 \qquad \sum_i\Theta_{y_i,x_i}=I,
\]
where the second equality is the adjoint of the first.

Direct expansion gives the exact polarization identity
\begin{equation}\label{eq:polarization}
 \Theta_{x+\eta y}-\Theta_x-N_{K/\Q}(\eta)\Theta_y
 =\overline\eta\,\Theta_{x,y}+\eta\,\Theta_{y,x}.
\end{equation}
Summing over $i$ and applying \eqref{eq:traceone} gives $I$ on the right. The left is an integer linear combination of rank-one operators of lattice vectors, since $N_{K/\Q}(\eta)\in\Z$. Combine repeated vectors to obtain
\begin{equation}\label{eq:signed-id}
 I=\sum_{v\in S'}z_v\Theta_v,\qquad z_v\in\Z.
\end{equation}
The prescribed set $S$ can be added to $S'$ with zero coefficients.
\end{proof}

The trace-one element removes the factor two that a purely real polarization argument would introduce. Projective dual generators replace a global integral basis; this is why no freeness assumption enters the proof.

\subsection{Step 2: a positive reserve containing prescribed vectors}

\begin{lemma}\label{lem:reserve}
For every finite $S'\subset A$, there are a positive integer $q$, a finite $T\supseteq S'$, and positive integers $c_v$ such that
\begin{equation}\label{eq:reserve}
 q I=\sum_{v\in T}c_v\Theta_v.
\end{equation}
\end{lemma}

\begin{proof}
Choose a sufficiently small positive rational $\epsilon$ so that
\[
 I-\epsilon\sum_{v\in S'}\Theta_v
\]
is positive definite. Choose a $K$-basis of $V$ inside $A$, with complex coordinate matrix $B$. Its Hermitian Gram matrix $H=B^*B$ has entries in $K$. If $u_v\in K^n$ is the coordinate column of $v$, then
\[
 I=BH^{-1}B^*,\qquad \Theta_v=Bu_vu_v^*B^*.
\]
The residual matrix
\[
 H^{-1}-\epsilon\sum_{v\in S'}u_vu_v^*
\]
is positive definite Hermitian over $K$. Hermitian $LDL^*$ decomposition has a triangular factor over $K$ and strictly positive diagonal entries in $K\cap\R=\Q$. This expresses the residual as a sum of positive rational multiples of rank-one matrices with vectors in $K^n$.

Multiply each vector by an integer so that its image in $V$ lies in $A$, and divide its coefficient by the corresponding square. Reinsert the vectors of $S'$ with coefficient $\epsilon$. Clearing all positive rational denominators gives \eqref{eq:reserve}. Zero vectors, if prescribed, may be retained with arbitrary positive coefficients. Combine duplicate vectors if necessary.
\end{proof}

\subsection{Step 3: absorb the signed coefficients}

\begin{proof}[Proof of Theorem~\ref{thm:projective}]
Use \eqref{eq:signed-id} and \eqref{eq:reserve}, extending $z_v=0$ outside the signed support. Define
\begin{equation}\label{eq:explicit-threshold}
 k_0=\max_{v\in T}
 \left\lceil\frac{R+(q-1)\max(0,-z_v)}{c_v}\right\rceil,
 \qquad M_0=qk_0.
\end{equation}
For $M\ge M_0$, write $M=kq+r$ with $0\le r<q$. Set
\[
 m_v=kc_v+rz_v.
\]
Then $k\ge k_0$ and every $m_v\ge R$. Moreover,
\[
 \sum_{v\in T}m_v\Theta_v
 =k\sum_{v\in T}c_v\Theta_v+r\sum_{v\in T}z_v\Theta_v
 =kqI+rI=MI.
\]
This proves the theorem with a fixed family independent of $M$.
\end{proof}

\begin{corollary}[Coprime and prescribed large bounds]
The corrected ideal $A=A_0\OO_E$ from Section~\ref{sec:unramified} has frames of every sufficiently large positive integer bound. In particular it has coprime pairs of bounds, and arbitrarily large bounds prime to any fixed integer. One can prescribe both old unit witnesses \eqref{eq:unit-witnesses} as frame members with any fixed lower multiplicity.
\end{corollary}

\begin{proof}
The ideal is Hermitian self-dual and projective, and $\OO_K$ contains the required trace-one element. Apply Theorem~\ref{thm:projective}. Consecutive sufficiently large integers give coprime bounds; integers congruent to one modulo a fixed integer give the other assertion.
\end{proof}

\subsection{A general finite local-unit witness principle}

The inherited two-element witness set is unusually convenient. A qualitative substitute exists for any fractional ideal.

\begin{lemma}[Finite simultaneous local generators]\label{lem:local-generators}
For a nonzero fractional ideal $A$ of a number field $E$, there is a finite set $S\subset A$ such that, for every rational prime $p$, some $y\in S$ satisfies
\[
 v_{\mathfrak P}(y)=v_{\mathfrak P}(A)
 \quad\text{for all }\mathfrak P\mid p.
\]
\end{lemma}

\begin{proof}
Choose $0\ne y_0\in A$. The integral ideal $(y_0)A^{-1}$ has finite support. Outside the finitely many rational primes below that support, $y_0$ already works. For each exceptional $p$, choose a nonzero class in every one-dimensional residue module $A/\mathfrak P A$, for $\mathfrak P\mid p$. The Chinese remainder theorem gives one element $y_p\in A$ representing those classes simultaneously. Its local valuation is exactly that of $A$ at every such prime. Take $S$ to consist of $y_0$ and these finitely many $y_p$.
\end{proof}

If $A$ is $\Gal(E/K)$-stable, one such element has the analogous simultaneous normalized-unit property under all $K$-embeddings, relative to local trivializations of $A$. A comparison with a particular global period normalization is still required. In our unramified benchmark the inherited source witnesses already satisfy that comparison as norm assertions.

\subsection{What this theorem does and does not settle}

For each fixed tower level, qualitative coprime frames and prescribed sufficiently large bounds are now available in the correct lattice. The earlier raw integral frame obstruction is therefore overcome by the fractional-lattice construction, not contradicted.

The proof deliberately records its threshold \eqref{eq:explicit-threshold}. It may be very large. Rational decomposition and denominator clearing do not by themselves bound $q$, $c_v$, $z_v$, the coordinates of the vectors, or $M_0$ uniformly in a growing tower. The source needs a small auxiliary bound $h\le b$ and exact final bound $M=p_0p_1$ in a prescribed window, with multiplicities at least $6fb$ and controlled conjugate norms. None of these quantitative inequalities follows from qualitative eventual existence. Establishing them is the first remaining substantial theorem.
